\documentclass[11pt]{article}
\usepackage[margin=1in]{geometry}
\usepackage{amsmath,amssymb,amsthm,graphicx}
\usepackage[hidelinks]{hyperref}
\newtheorem{theorem}{Theorem}
\newtheorem{lemma}{Lemma}
\newcommand{\Cb}{\mathcal C_b}
\setlength{\parindent}{0pt}
\setlength{\parskip}{0.5em}

\title{A Sublinear Two-Sided Near-Additivity Counterexample}
\author{Run \texttt{fa\_banach\_001}, lane 6}
\date{22 July 2026}

\begin{document}
\maketitle

\textbf{Status.} Candidate full counterexample, likely valid, for human
review.

\section{Source question}

Brice\~no and Iommi prove that if a sequence
$(f_n)_n\subseteq\Cb(X)$ satisfies
\begin{equation}\label{eq:source-defect}
 \bigl\|f_{n+m}-f_n-f_m\circ T^n\bigr\|_\infty\leq C_n+C_m
 \qquad(m,n\in\mathbb N)
\end{equation}
for a nondecreasing nonnegative sequence such that
$\sum_{n\geq1}C_n/n^2<\infty$, then $(f_n)_n$ is asymptotically additive
\cite[Proposition~5.1]{BI}. Their Question~5.2 asks whether the summability
assumption can be replaced by only
\[
 0\leq C_n\leq C_{n+1},\qquad \frac{C_n}{n}\longrightarrow0.
\]
The question occurs on PDF page~27.

\begin{center}
\includegraphics[width=\textwidth]{figures/open_problem_crop.png}
\end{center}

\section{Result}

\begin{theorem}\label{thm:main}
The proposed weakening in Question~5.2 is false. More precisely, there are a
compact dynamical system $(X,T)$, a sequence $(f_n)_n\subseteq\Cb(X)$, and a
sequence $(C_n)_{n\geq0}$ such that:
\begin{enumerate}
\item $C_0=0$, and $(C_n)_{n\geq1}$ is positive, nondecreasing, and even
subadditive;
\item $C_n/n\to0$;
\item \eqref{eq:source-defect} holds for all $m,n\geq0$;
\item $(f_n)_n$ is not asymptotically additive.
\end{enumerate}
In fact, the failure in the last item has normalized distance exactly one
from the additive sequences.
\end{theorem}

\section{Proof intuition}

On a one-point system, asymptotic additivity is simply convergence of
$f_n/n$. We therefore want a bounded, nonconvergent slope $h(n)=f_n/n$ which
changes so slowly that $f_n=nh(n)$ is additive up to a sublinear error. The
choice
\[
 h(x)=\sin(\log\log(x+e^4))
\]
does exactly this: its phase tends to infinity, so it keeps oscillating, while
its derivative is bounded by
$((x+e^4)\log(x+e^4))^{-1}$. For two comparable summands this derivative
directly gives an $O(n/\log n)$ defect. If one summand is much smaller, its
coefficient absorbs the larger change in phase. Splitting at $m=\sqrt n$
makes this uniform and yields the explicit error
$C_n=15n/\log(n+e^4)$.

\section{Proof}

Put $B=e^4$ and define, for $x\geq0$,
\[
 q(x)=\log\log(x+B),\qquad h(x)=\sin q(x).
\]
Then
\begin{equation}\label{eq:derivative}
 |h'(x)|\leq q'(x)=\frac{1}{(x+B)\log(x+B)}.
\end{equation}

Take $X=\{*\}$ with $T(*)=*$. Set
\[
 f_0(*)=0,\qquad f_n(*)=n h(n)\quad(n\geq1),
\]
and
\begin{equation}\label{eq:C}
 C_0=0,\qquad C_n=\frac{15n}{\log(n+B)}\quad(n\geq1).
\end{equation}
Every $f_n$ is of course bounded and continuous.

We first verify the error estimate. By symmetry it is enough to consider
$1\leq m\leq n$. Algebra and \eqref{eq:derivative} give
\begin{align}
 |f_{n+m}-f_n-f_m|
 &\leq n|h(n+m)-h(n)|+m|h(n+m)-h(m)| \notag\\
 &\leq n\int_n^{n+m}\frac{dx}{(x+B)\log(x+B)}
   +m\int_m^{n+m}\frac{dx}{(x+B)\log(x+B)}.\label{eq:split}
\end{align}
The first integral term is at most
\begin{equation}\label{eq:first}
 \frac{nm}{(n+B)\log(n+B)}\leq\frac{n}{\log(n+B)}.
\end{equation}
Denote the second term by $E_{m,n}$. We claim that
\begin{equation}\label{eq:second}
 E_{m,n}\leq\frac{14n}{\log(n+B)}.
\end{equation}

Suppose first that $m\leq\sqrt n$. Since $m\leq n$ and $\log B=4$,
\begin{align*}
 E_{m,n}
 &=m\log\!\left(\frac{\log(n+m+B)}{\log(m+B)}\right)\\
 &\leq \frac{m}{4}\log(2(n+B)).
\end{align*}
Here we used $\log y\leq y$ for $y>0$. Since
$n+B\leq(B+1)n<e^5n$ and $2(n+B)<e^6n$, while
$\log n\leq2n^{1/4}$ for $n\geq1$, it follows that
\begin{align*}
 \frac{E_{m,n}\log(n+B)}{n}
 &\leq \frac{\log(n+B)\log(2(n+B))}{4\sqrt n}\\
 &\leq\frac{(5+2n^{1/4})(6+2n^{1/4})}{4\sqrt n}\\
 &=\frac{30}{4\sqrt n}+\frac{22}{4n^{1/4}}+1\leq14.
\end{align*}

Now suppose that $m>\sqrt n$. Using $\log(1+u)\leq u$ twice,
\begin{align*}
 E_{m,n}
 &=m\log\!\left(
  1+\frac{\log(1+n/(m+B))}{\log(m+B)}\right)\\
 &\leq\frac{m}{\log(m+B)}\log\!\left(1+\frac{n}{m+B}\right)
 \leq\frac{n}{\log(m+B)}.
\end{align*}
Moreover $(m+B)^2>n+B$, so
$\log(n+B)<2\log(m+B)$. Consequently
\[
 E_{m,n}\leq\frac{2n}{\log(n+B)},
\]
which is stronger than \eqref{eq:second}. Combining
\eqref{eq:split}--\eqref{eq:second} yields
\[
 |f_{n+m}-f_n-f_m|
 \leq\frac{15n}{\log(n+B)}=C_n\leq C_n+C_m.
\]
The cases $m=0$ or $n=0$ are identities, so \eqref{eq:source-defect} holds
for all nonnegative $m,n$.

Next consider the hypotheses on the errors. The function
$x\mapsto x/\log(x+B)$ is strictly increasing on $[0,\infty)$, because its
derivative has the sign of
\[
 \log(x+B)-\frac{x}{x+B}>4-1>0.
\]
Thus $(C_n)$ is nondecreasing. Also
\[
 \frac{C_n}{n}=\frac{15}{\log(n+B)}\longrightarrow0.
\]
It is even subadditive: since $x\mapsto1/\log(x+B)$ decreases,
\begin{align*}
 \frac{n+m}{\log(n+m+B)}
 &\leq\frac{n}{\log(n+B)}+\frac{m}{\log(m+B)}.
\end{align*}
For comparison with Proposition~5.1, the excluded summability condition does
indeed fail, since
\[
 \sum_{n\geq1}\frac{C_n}{n^2}
 =15\sum_{n\geq1}\frac{1}{n\log(n+B)}=\infty.
\]

It remains to show failure of the conclusion. Choose integers $n_k^+$ and
$n_k^-$ nearest, respectively, to
\[
 \exp\!\bigl(\exp(\tfrac\pi2+2\pi k)\bigr)-B,
 \qquad
 \exp\!\bigl(\exp(\tfrac{3\pi}2+2\pi k)\bigr)-B.
\]
Since $q'(x)\to0$, rounding changes $q$ by $o(1)$. Hence
\[
 h(n_k^+)\longrightarrow1,\qquad h(n_k^-)\longrightarrow-1.
\]
Thus $f_n/n=h(n)$ does not converge.

More exactly, every $g\in\Cb(X)$ is a scalar $\lambda$, and
$S_ng=n\lambda$. Therefore
\[
 \limsup_{n\to\infty}\frac{\|f_n-S_ng\|_\infty}{n}
 =\limsup_{n\to\infty}|h(n)-\lambda|
 \geq\max\{|1-\lambda|,|-1-\lambda|\}\geq1.
\]
Equality is attained at $\lambda=0$. Hence the infimum in the definition of
asymptotic additivity is exactly one, rather than zero. This proves
Theorem~\ref{thm:main}. \qed

\section{Verification, scope, and novelty}

The proof uses no external theorem beyond the source definitions. The included
finite-range verifier exhaustively checks the defect for all
$1\leq m\leq n\leq800$, checks selected logarithmic-scale pairs up to
$n=10^{12}$, and checks finite-range monotonicity and subadditivity of
$(C_n)$. These computations are not used in the proof.

The bounded novelty search on 22 July 2026 covered the run's four lightweight
indexes, the exact arXiv id and question wording, and web/arXiv searches for
the phrases \texttt{C\_n + C\_m}, \texttt{asymptotically additive}, and
\texttt{sublinear nearly additive}, with close variants. No separate paper
explicitly resolving Question~5.2 was found. The source notes only that an
analogous one-sided nearly-subadditive scalar question was answered negatively
by F\"uredi. This packet does not rely on that construction. Novelty is not
certified beyond the bounded search.

The counterexample settles the yes/no weakening in Question~5.2. It does not
seek the sharp additional regularity on $(C_n)$ that would restore the
conclusion. Human review should focus on the two cases in
\eqref{eq:second}; the one-point reduction and the final oscillation are
otherwise direct.

\begin{thebibliography}{9}
\bibitem{BI}
R.~Brice\~no and G.~Iommi,
\emph{Ergodic theorems for set maps under weak forms of additivity},
arXiv:2501.14581v1 (2025).
\end{thebibliography}

\end{document}
